% ============================================================
% PAGE 2 — Introduction
% ============================================================

\section{Introduction}
\label{sec:intro}

Object detection—the simultaneous localization and semantic
classification of every instance of interest within an
image—occupies a foundational role in modern computer vision.
Autonomous vehicles must identify pedestrians, cyclists, and
road signs within millisecond-scale budgets; medical imaging
systems must locate lesions and anatomical landmarks with
sub-pixel accuracy; and surveillance platforms must track
multiple objects across complex, dynamic scenes. In each
setting, the tension between \textit{speed} and
\textit{accuracy} is real and consequential.

Over the past decade, deep learning has transformed the detection
landscape. Two-stage detectors such as Faster R-CNN
\cite{ren2015faster} decompose the problem into a region-proposal
step followed by per-proposal classification and regression, yielding
high accuracy at the cost of elevated computational budgets.
One-stage detectors---YOLO \cite{redmon2016yolo} and its
descendants---collapse these stages into a single forward pass,
achieving competitive accuracy at substantially lower latency.
Yet both families share a common architectural assumption: a
\textit{single backbone} extracts all features, and the subsequent
neck and head merely reorganize and decode those features.

This assumption has come under scrutiny in the transformer era.
Vision Transformers (ViTs) \cite{dosovitskiy2021vit}, built on
the multi-head self-attention primitive of Vaswani
\textit{et al.}\ \cite{vaswani2017attention}, can model
relationships between image patches separated by hundreds of pixels---
a capability that convolution's local receptive fields cannot
replicate without extreme depth. End-to-end detectors such as DETR
\cite{carion2020detr} and its hierarchical extension Swin-Transformer
\cite{liu2021swin} have accordingly achieved state-of-the-art results
on standard benchmarks. The cost, however, is considerable: full
self-attention scales quadratically with the number of tokens, making
high-resolution detection expensive.

\subsection{Motivation}
\label{sec:motivation}

The core insight motivating FusionDet is that CNNs and transformers
are \textit{complementary} rather than competing. A CNN processes
every spatial location with learnable convolutional filters: it is
parameter-efficient, hardware-friendly, and naturally equivariant to
local translations. A transformer processes every token in relation
to every other token: it is globally aware, dynamically reweights
distant context, and is less constrained by the locality of
convolutional kernels. Object detection benefits from both
properties---local texture and edge information is critical for
precise localization, while global context is critical for resolving
ambiguous classifications in cluttered scenes.

Prior work has combined CNNs and transformers in various ways:
inserting transformer blocks at the bottleneck of a CNN backbone,
replacing convolutional layers with depthwise attention, or using
a transformer decoder to query a CNN-derived feature map. Each of
these hybrids privileges one paradigm over the other, using the
secondary modality as a post-hoc enhancement rather than as an equal
partner. FusionDet, by contrast, operates \textit{dual backbones in
parallel} and fuses their outputs at every scale level through a
learned Adaptive Feature Fusion module, ensuring that both
representational streams contribute symmetrically to every
downstream computation.

\subsection{Contributions}
\label{sec:contributions}

This paper makes the following four original contributions:

\begin{enumerate}
  \item \textbf{Dual-Backbone Architecture.} We present a parallel
    CNN--transformer backbone that processes each input image
    simultaneously through a CSP-Darknet branch and a linear-attention
    transformer branch, producing spatially and semantically aligned
    feature tensors at four resolution levels.

  \item \textbf{Adaptive Feature Fusion (AFF) Module.} We propose a
    trainable gating module that learns, \textit{per-channel and
    per-spatial-location}, how to blend the CNN and transformer
    outputs---going beyond the fixed weighted sums used in prior
    fusion work.

  \item \textbf{Bidirectional Path Aggregation Neck with Lightweight
    Channel Attention.} We integrate a LCA module into every level of
    a bidirectional PANet neck, providing dynamic channel
    recalibration at negligible parameter and compute overhead.

  \item \textbf{Comprehensive Empirical Evaluation.} We present
    thorough ablation studies on MS-COCO that isolate the contribution
    of each proposed component, and we compare FusionDet against
    twelve contemporary detectors across five speed--accuracy
    operating points.
\end{enumerate}

The remainder of the paper is organized as follows.
Section~\ref{sec:related} surveys related work.
Section~\ref{sec:problem} formalizes the detection problem and
introduces notation. Section~\ref{sec:arch} presents the full
FusionDet architecture. Section~\ref{sec:results} reports
experimental results, and Section~\ref{sec:conclusion} concludes.
